\documentclass[11pt]{article}
\usepackage[utf8]{inputenc}
\usepackage[T1]{fontenc}
\usepackage[spanish]{babel}
\usepackage{makeidx}
\usepackage[backend=biber,style=authoryear]{biblatex}
\addbibresource{refs.bib}
\makeindex

\begin{document}

\section{Fuentes}

Los filtros de Butterworth\index{Butterworth} se describen en \parencite{butterworth1930}, y la aproximación de Chebyshev\index{Chebyshev} aparece tratada en \textcite{papoulis1958}. Un párrafo largo en una sola línea que repite términos indexados, como Butterworth\index{Butterworth} y Cauer\index{Cauer}, para que el índice alfabético tenga entradas con más de una página y haya trabajo para makeindex y para biber en la misma corrida.

\nocite{cauer1958}

\printbibliography

\printindex

\end{document}
